% !PRINTED 577
\weberpage{577}
The form of this development can be derived from the product formulae [\S\ 24, (11)]:
\[
\lambda = q^{-\frac{1}{2}} \Pi\left(\frac{1+q^{2\nu-1}}{1+q^{2\nu}}\right)^4 + 16\,q^{\frac{1}{2}}\ \Pi\left(\frac{1+q^{2\nu}}{1+q^{2\nu-1}}\right)^4
\]
and one finds:
\weq{3}{\lambda = q^{-\frac{1}{2}}(1 + \lambda_1 q + \lambda_2 q^2 + \lambda_3 q^3 + \cdots),}
in which the constants $\lambda_1, \lambda_2, \lambda_3, \ldots$ are \gesperrt{integral rational numbers}.

If we further set
\weq{4}{x = \sqrt{\varkappa}\ \operatorname{sn}\left(\frac{w}{\sqrt{\varkappa}},\varkappa\right),}
then, by \S\ 45, (9), (4) remains unchanged when $\varkappa$ is interchanged with $1:\varkappa$. Now $x$ satisfies the differential equation
\weq{5}{\left(\frac{dx}{dw}\right)^2 = 1 - \tfrac{1}{4}\lambda x^2 + x^4.}
If, then, one develops $x$ by \gesperrt{Taylor's} theorem into a series in ascending powers of $w$:
\weq{6}{x = w + X_1 w^3 + X_2 w^5 \ldots,}
then the coefficients $X_1, X_2, \ldots$ are \gesperrt{integral rational functions of $\lambda$ with rational numerical coefficients}, as follows from (5) by the method of indeterminate coefficients.

If one uses the notation of \S\ 42, then
\weq{7}{x = \sqrt{\varkappa}\ \operatorname{sn} v = \frac{\vartheta_{11}(u)}{\vartheta_{01}(u)},}
and, by \S\ 57, IV., with a somewhat modified meaning of the $A, B, C, D$:
\weq{8}{\begin{aligned}
\sqrt{\varkappa}\ \operatorname{sn} mv &= \frac{xA(x)}{D(x)}, \\
\operatorname{cn} mv &= \sqrt{1-\frac{x^2}{\varkappa}}\ \frac{B(x)}{D(x)}, \\
\operatorname{dn} mv &= \sqrt{1-\varkappa x^2}\ \frac{C(x)}{D(x)}.
\end{aligned}}
Here $A(x), B(x), C(x), D(x)$ are integral rational functions of $x$ of degree $\frac{1}{2}(m^2-1)$, and it is to be seen from the recursion formulae \S\ 57, (13), (14) that the functions $A, D$ and the product $BC$ are integral rational functions of $\lambda$, whose
% [source sic: the printed degree (m^2-1)/2 conflicts with the m^2-1 roots and later degree formulae]
%CONT>>
